\section{Глава~\ref{sec:gaba}. \nt{ГЛУТ} и \nt{ГАМК} вместе}

Логические и динамические утверждения главы выведены в ней самой из линейного анализа и закона Ома; опыт показывает, что схемы, на которые они опираются, --- запрет с задержкой, прямое торможение вслед за возбуждением, стабилизация торможением, ритм петли \nt{ГЛУТ}--\nt{ГАМК}, --- в мозге действительно работают.

{\footnotesize
\setlength{\tabcolsep}{3pt}\renewcommand{\arraystretch}{1.15}
\begin{longtable}{>{\raggedright}p{0.5cm}>{\raggedright}p{4.0cm}>{\raggedright}p{5.6cm}>{\raggedright}p{2.3cm}>{\raggedright\arraybackslash}p{1.7cm}}
\toprule
№ & Утверждение & Опыт и что измерено & Источник & Статус \\
\midrule\endfirsthead
\toprule
№ & Утверждение & Опыт и что измерено & Источник & Статус \\
\midrule\endhead
1 & \nt{ГАМК}-нейронов в коре около пятой части & Иммуноокраска на ГАМК в $10$ областях коры обезьяны: около $25\,\%$ нейронов в $8$ областях, меньше $20\,\%$ в первичной зрительной коре. Обзор разнообразия тормозящих нейронов коры & \cite{Hendry1987,Markram2004} & измерено \\
2 & Что делает торможение, во многом решает место посадки: дендрит, тело, место рождения импульса, окончание чужого провода & Обзор: классы \nt{ГАМК}-нейронов коры различаются мишенью на получателе. Одновременная запись с тела и дендрита пирамидного нейрона: прямое торможение гораздо сильнее на теле, чем на дендритах. Торможение на окончании чужого провода в спинном мозге уменьшает выброс медиатора одной входной линии (обзор измерений) & \cite{Markram2004,Pouille2001,Rudomin1999} & измерено \\
3 & Смена знака через посредника стоит одной задержки, около миллисекунды; торможение, пришедшее вслед за возбуждением, сужает окно совпадения & Пирамидные нейроны гиппокампа крысы в срезах: торможение через одного посредника приходит вслед за прямым возбуждением с короткой задержкой, и окно, в котором два возбуждающих входа суммируются до порога, меньше $2\ms$. На дендритах, где это торможение слабее, окно шире & \cite{Pouille2001} & измерено \\
4 & Запрет $a\wedge\bar b$~\eqref{eq:veto} с ИЛИ и постоянной единицей функционально полон (утверждение~\ref{prop:gabacomplete}) & Доказательство в главе. Классический результат: сеть пороговых элементов с тормозящими входами реализует любое логическое выражение. Утверждение о модели, опыта нет & \cite{McCullochPitts1943} & теория \\
5 & Запрет с задержкой даёт детектор направления движения с оптимальной скоростью $v^*=\Delta x/d$~\eqref{eq:vstar} & Сетчатка кролика: избирательность к направлению создаётся торможением, которое при движении в «нулевом» направлении гасит возбуждение. Раздельная запись возбуждающего и тормозящего входов избирательных клеток показала, что звёздчатые амакриновые клетки (\nt{ГАМК}) тормозят их несимметрично --- только отростками, направленными в «нулевую» сторону. Формула $v^*$ --- вывод главы & \cite{Barlow1965,Fried2002} & измерено \\
6 & Шунтирующее торможение делит напряжение, а не вычитает~\eqref{eq:shunt} (рис.~\ref{fig:shunt}) & Закон Ома для делителя из трёх проводимостей при равновесном потенциале хлора около покоя; для нейрона без импульсов & вывод в главе & теория \\
7 & На частоту импульсов шунт действует вычитанием, а деление получается от торможения вместе с уравновешенным фоновым шумом & Модель: у срабатывающего нейрона ток через шунт почти постоянен, и эффект на частоту вычитающий. Опыт: в пирамидные нейроны соматосенсорной коры крысы (срезы) через динамический фиксатор вводили фоновые возбуждающие и тормозящие проводимости; одновременный уравновешенный рост обеих делит крутизну ответа без заметного сдвига частоты. Обзор: деление на общую активность встречается во многих отделах мозга & \cite{HoltKoch1997,Chance2002,Carandini2012} & измерено \\
8 & При $\beta>1$ взаимное торможение выбирает одного победителя (утверждение~\ref{prop:wta}, \eqref{eq:wta}) & Собственные числа $-(1\pm\beta)/\tau$ --- вывод в главе. Частоты $0.73$ и $0.53$ при $\beta=0.5$ --- неподвижная точка $r_{1,2}=(I_{1,2}-\beta I_{2,1})/(1-\beta^2)$; скрипта в \texttt{models/} нет. Прямого опыта в главе не приводится & вывод в главе & теория \\
9 & Сброс памяти сигналом через \nt{ГАМК}; две группы с взаимным торможением --- защёлка & Следует из рис.~\ref{fig:bistable} и утверждения~\ref{prop:wta}. Опыта нет & вывод в главе & теория \\
10 & Равновесие петли \nt{ГЛУТ}--\nt{ГАМК} устойчиво, если торможение достаточно сильное и достаточно быстрое (утверждение~\ref{prop:ei}) & Модель двух популяций~\eqref{eq:ei}; условия (i) и (ii) --- определитель и след её матрицы, выведены в главе & \cite{WilsonCowan1972} & теория \\
11 & При $w_{EE}=1.5$ и $w_{EI}w_{IE}=2$ быстрое торможение ($\tau_I=2\ms$) держит $E^*=0.67h$, медленное ($30\ms$) раскачивает колебания (рис.~\ref{fig:ei}) & Численное решение~\eqref{eq:ei}, скрипт \texttt{figures/make\_ei.py} & расчёт & модель \\
12 & Кора работает в режиме стабилизации торможением; подпись --- подтолкнёшь \nt{ГАМК}-нейроны, и их частота упадёт~\eqref{eq:isn} & Теория предсказала парадоксальный ответ. Опыт: внутриклеточные записи в первичной зрительной коре кошки --- при подавлении ответа раздражителем вокруг тормозящая проводимость не растёт, а падает вместе с возбуждающей. Оптогенетическое возбуждение PV-нейронов мыши в зрительной, соматосенсорной и моторной коре снижает частоту тормозящих нейронов, в том числе без раздражителя и при лёгком наркозе. Коэффициент $-1/3$ при числах рис.~\ref{fig:ei} --- вывод главы & \cite{Tsodyks1997isn,Ozeki2009,Sanzeni2020} & измерено \\
13 & Петля «\nt{ГЛУТ} возбуждает \nt{ГАМК}, \nt{ГАМК} тормозит \nt{ГЛУТ}» порождает быстрый ритм & Оптогенетическая стимуляция быстрых \nt{ГАМК}-нейронов в соматосенсорной коре мыши на частотах $8$--$200$\,Гц избирательно усиливает гамма-ритм ($20$--$80$\,Гц, период $12$--$50\ms$); стимуляция \nt{ГЛУТ}-нейронов усиливает только более медленные ритмы. В главе период назван $10$--$30\ms$, то есть верхняя часть этой полосы & \cite{Cardin2009,Buzsaki2012} & измерено \\
\bottomrule
\end{longtable}}
